\section{Cómo pueden los investigadores jóvenes entrar a la investigación en arquitecturas}

Una vez expuestas la trayectoria técnica de Attention y las restricciones de hardware, esta sección regresa la discusión al método de investigación. La investigación en arquitecturas no consiste solo en leer artículos ni únicamente en saber escribir fórmulas; requiere estar cerca de modelos reales, cómputo real y restricciones de hardware reales. Para un investigador joven, la pregunta más importante es cómo entrar a un entorno donde realmente pueda verificarse el comportamiento a escala.

El consejo para los investigadores jóvenes al final del programa es directo: hacer una pasantía en un lab y seguirle el paso a la frontier, porque trabajar en arquitecturas exige cómputo. El consejo suena sencillo, pero tiene razones muy concretas: la investigación en arquitecturas difícilmente puede evaluarse solo con experimentos a pequeña escala, y una idea que funciona en un modelo pequeño no necesariamente funciona en un modelo grande.

\begin{figure}[H]
\centering
\includegraphics[width=0.96\textwidth]{young-researcher-advice.png}
\caption{El camino para entrar a la investigación en arquitecturas: arqueología de la literatura, cómputo, un frontier lab y conciencia del hardware son todos indispensables. Diagrama conceptual propio, elaborado a partir de la conversación en 01:42:23--01:43:26.}
\end{figure}

\begin{knowledgebox}{Lectura de la figura: la investigación en arquitecturas necesita estar cerca de las restricciones reales}
Leer artículos antiguos para construir un mapa, encontrar problemas útiles, obtener suficiente cómputo para verificar el comportamiento a escala, entrar a un frontier lab para acercarse a un entorno de entrenamiento real y, después, usar implementaciones de código abierto para que la técnica se difunda. Es un camino que va de la comprensión a la verificación y de ahí al impacto.
\end{knowledgebox}

\subsection{Por qué el cómputo es el boleto de entrada}

Muchos problemas de la investigación en arquitecturas solo se manifiestan a partir de cierta escala. La expresividad, la estabilidad del entrenamiento, el KV cache, el decoding, los kernels, el tamaño de lote y el rendimiento del hardware no pueden verificarse por completo con pequeños experimentos en una sola máquina. Sin cómputo, es muy difícil saber si una arquitectura es una línea de desarrollo con futuro o una ilusión de la escala pequeña.

\subsection{Resumen de la sección}

Entrar a la investigación en arquitecturas exige tanto criterio algorítmico como conciencia de ingeniería y de hardware. Un investigador joven no puede limitarse a leer los artículos más recientes ni a hacer experimentos de juguete; lo más importante es acercarse a las restricciones reales de la frontier.
